\subsection{The extension threshold of the eighth character}
\label{subsec:eta-eight-extension-threshold}

Complex conjugation gives an arithmetic obstruction to extending the
new character through the preceding normalizations.

\begin{proposition}\label{prop:eta-eight-no-extension-d3}
The character $\eta_8$ on $D_7$ extends to the invariant character
\[
 \alpha_8:D_5\longrightarrow\mathbf F_2,\qquad
 \alpha_8(1,f)=[S^4TST^2]f\pmod2,
\]
but has no $G$-invariant extension to $D_3$. More precisely, no
nonzero linear combination of $\eta_7$ and $\alpha_8$ on $D_5$
extends to a $G$-invariant character of $D_3$.
\end{proposition}


\begin{proof}
The exact weight-seven theorem puts the degree-seven part of every
$D_5$ word in $8\mathbf Z_2\mathcal L_7$, where $\mathcal L_7$
is the primitive vector of \eqref{eq:weight-seven-primitive-vector}.
Thus its Magnus expansion over $\mathbf F_2$ starts in degree eight.
The first nonzero term is additive under the normalized GT product,
and conjugation acts by $\lambda_2^8$, which is one modulo two.
This follows by substituting the linear terms in the GT automorphism:
higher terms raise degree, and a homogeneous ambiguity commuting with
$T$ would be a pure power of $T$, absent from a commutator word.
Hence $\alpha_8$ is an invariant character and restricts to $\eta_8$.

On $M=\ker\lambda$ use the additive coordinates
\[
 a=m_3,\qquad T_5=12[S^4T]f-11a,\qquad
 T_7=360[S^6T]f-1020[S^4T]f+661a.
\]
Their conjugation weights are $\lambda_2^3,\lambda_2^5,\lambda_2^7$.
Together with the rational Kummer homomorphisms they imply
\[
 [M,M]\subset D_7,\qquad [M,D_7]\subset D_8.
\]
For the first inclusion all four families vanish on commutators;
then $T_5=12m_5$ and $T_7=360m_7$, and $\mathbf Z_2$ has no torsion.
For the second inclusion use the weight-eight law on $D_7$ and
$\lambda(M)=1$. Therefore $M/D_8$ has nilpotency class at most two,
and its commutator subgroup is contained in $D_7/D_8\simeq\mathbf Z_2$.
The previously proved commutator formula is
\begin{equation}\label{eq:eta-eight-commutator-pairing}
 m_8([z,z'])=\frac7{24}\bigl(a(z)T_5(z')-a(z')T_5(z)\bigr)
 \qquad(z,z'\in M).
\end{equation}

Let $c\in H$ be complex conjugation, and choose $h\in H$ with
$\lambda_2(h)=3$. Put $t=\lambda(h)$ and $\sigma=[h,c]$.
Since $c^2=1$, $\lambda(c)=-1$ and $\rho_2(c)=0$, one has
\[
 c\sigma c^{-1}=\sigma^{-1},\qquad
 \rho_2(\sigma)=2\rho_2(h).
\]
The established parity law makes $\rho_2(h)$ odd at two, and
$\rho_2(\sigma)\equiv-2m_3(\sigma)\pmod4$ makes
\[
 A=a(\sigma),\qquad B=T_5(\sigma),\qquad C=T_7(\sigma)
\]
all odd. The last two assertions also follow directly from their
formulas above. The odd-prime Kummer coordinates of $\sigma$ need
not vanish.

For $e\in\widehat{\mathbf Z}$ put
$D_e(z)=hzh^{-1}z^{-e}$, and define
\[
 z_1=D_t(\sigma),\qquad z_3=D_{t^3}(z_1),\qquad
 z_5=D_{t^5}(z_3).
\]
The weight laws give $z_1\in D\cap H$, $z_3\in D_3\cap H$ and
$z_5\in D_5\cap H$. Moreover
\[
 m_7(z_5)=\frac{(3^7-3)(3^7-3^3)(3^7-3^5)}{360}C
        =128\cdot199017\,C,
\]
so $\eta_7(z_5)=1$. Every invariant quadratic character on $D_3$
kills $z_5$, since conjugation fixes its value on $z_3$ and
$t^5$ is odd at two.

A second arithmetic element is $v=cz_3c^{-1}z_3$.
Every rational Kummer coordinate and each of $a,T_5,T_7$ changes
sign under $c$, so $v\in D_7$. Every invariant character on $D_3$
also kills $v$. We compute $m_8(v)$ in the class-two quotient.
Put $X=\sigma$, $Y=h\sigma h^{-1}$, $Z=h^2\sigma h^{-2}$.
The identity $chc^{-1}=X^{-1}h$ gives
\[
 cXc^{-1}=X^{-1},\quad
 cYc^{-1}=X^{-1}Y^{-1}X,\quad
 cZc^{-1}=X^{-1}Y^{-1}Z^{-1}YX.
\]
For profinite exponents $p,r$, set $W=ZY^{-p}(YX^{-p})^{-r}$.
Class-two collection gives
\begin{equation}\label{eq:eta-eight-conjugation-norm}
 (cWc^{-1})W=[X,Y]^{pr(p-1)-p-r}
                  [X,Z]^{1-pr}[Y,Z]^{1+p+r}.
\end{equation}
Here is a direct verification. Encode the normal form
$X^aY^bZ^c[X,Y]^d[X,Z]^e[Y,Z]^f$ by $(a,b,c,d,e,f)$.
Multiplication is
\[
 (a,b,c,d,e,f)(A,B,C,D,E,F)
 =(a+A,b+B,c+C,d+D-Ab,e+E-Ac,f+F-Bc),
\]
and its $n$th power has coordinates
\[
 (na,nb,nc,nd-\tbinom n2ab,
             ne-\tbinom n2ac,nf-\tbinom n2bc).
\]
The normal form of $W$ is
\[
 (pr,-p-r,1,\ p^2r+pr(r-1)/2,-pr,p+r).
\]
Conjugation by $c$ sends a coordinate tuple to
$(-a,-b,-c,d+b,e+c,f+c)$. Multiplying these two tuples gives
\eqref{eq:eta-eight-conjugation-norm}. Integer exponent identities
extend to profinite exponents by continuity, including the integral
binomial polynomial $n(n-1)/2$.

For $W=z_3$, the exponents are $p=t,r=t^3$. Only after applying
$m_8$ to the central commutator group do their two-adic values become
$3,27$. Thus the exponents in \eqref{eq:eta-eight-conjugation-norm}
are $132,-80,31$ in that coordinate. Formula
\eqref{eq:eta-eight-commutator-pairing} gives
\[
 m_8([X,Y])=63AB,\quad m_8([X,Z])=17010AB,\quad
 m_8([Y,Z])=413343AB.
\]
Consequently
\[
 m_8(v)=(132\cdot63-80\cdot17010+31\cdot413343)AB
       =11461149\,AB
\]
is odd. Hence $(\eta_7(v),\alpha_8(v))=(0,1)$.

Suppose an invariant character on $D_3$ restricts to
$\mu\eta_7+\nu\alpha_8$ on $D_5$. Evaluation at $v$ forces
$\nu=0$; evaluation at $z_5$ then forces $\mu=0$.
The value $\alpha_8(z_5)$ is unnecessary. Finally, an extension
of $\eta_8$ to $D_3$ would restrict on $D_5$ to
$\alpha_8+\epsilon\eta_7$, because $D_5/D_7\simeq128\mathbf Z_2$.
Both possibilities have been excluded.
\end{proof}

\begin{corollary}\label{cor:d3-character-kernel-d8}
One has
\[
 \ker\!\left(H^1(D_3,\mathbf F_2)^G
                    \longrightarrow H^1(D_8,\mathbf F_2)\right)
       =\mathbf F_2\eta_5.
\]
\end{corollary}
\begin{proof}
Let $\chi$ be an invariant character on $D_3$ that vanishes on $D_8$.
Since $D_7/D_8\simeq\mathbf Z_2$ through $m_8$, its restriction
to $D_7$ is zero or $\eta_8$. Subtracting the corresponding multiple
of $\alpha_8$ from $\chi|_{D_5}$ makes it vanish on $D_7$.
It then factors through $D_5/D_7\simeq128\mathbf Z_2$, so
$\chi|_{D_5}=\mu\eta_7+\nu\alpha_8$.
The proposition forces both coefficients to be zero.
Thus $\chi$ factors through $D_3/D_5\simeq32\mathbf Z_2$ and is
a multiple of $\eta_5$. Conversely $\eta_5$ is invariant on $D_3$
and vanishes on $D_5$, hence on $D_8$.
\end{proof}
